Expert witnesses in redistricting litigation face a recurring methodological
challenge: opposing counsel argues that any single compactness metric was
chosen because it favours the expert's side, and that a different metric
would tell a different story.
This ``metric shopping'' objection is legitimate: as documented in K.0,
the six standard compactness metrics have pairwise correlations ranging from
$r = 0.27$ (PWC--CH) to $r = 0.79$ (Reock--CH), confirming that different
metrics can and do give different answers about the same district.

The solution is a composite profile: present all six metrics, explain
what each measures, and let the court weigh them.
A composite profile is robust to metric-shopping objections because the expert
is not selecting the most favourable metric---they are presenting all metrics
and explaining the full picture.
This approach has been used in academic redistricting literature~\citep{niemi1990measuring}
and is increasingly standard in expert reports filed in state redistricting litigation.

\subsection{The bisect Platform}

The \textsc{bisect} platform supports composite reporting natively.
The command:
\begin{verbatim}
bisect label-analyze <label> --year 2020 --types all
\end{verbatim}
computes all six K-track metrics for every district in a single pass,
producing structured JSON output suitable for inclusion in expert reports.
The JSON is self-documenting with metric names, values, ranges, and projection
specifications; formal certification as a court-filing appendix requires counsel
review and expert witness attestation.

\subsection{Structure of This Paper}

Section~2 surveys court cases where each metric has been cited.
Section~3 explains the composite profile methodology.
Section~4 addresses weighting when metrics disagree.
Section~5 documents the \texttt{bisect label-analyze} output format.
Section~6 provides template court-filing language.
Section~7 reports the simulated panel agreement analysis.
